\documentclass[12pt, titlepage]{article}

\usepackage{fullpage}
\usepackage[numbers]{natbib}
\usepackage{booktabs}
\usepackage{tabularx}
\usepackage{graphicx}
\usepackage{float}
\usepackage[margin=1in]{geometry}
\usepackage{amsmath}
\usepackage{amssymb}
\usepackage{listings}
\usepackage{xcolor}
\usepackage{hyperref}
\hypersetup{
    colorlinks,
    citecolor=blue,
    filecolor=black,
    linkcolor=red,
    urlcolor=blue
}

\lstset{
  language=C,
  basicstyle=\ttfamily\small,
  keywordstyle=\color{blue}\bfseries,
  commentstyle=\color{gray},
  stringstyle=\color{orange},
  breaklines=true,
  frame=single,
  numbers=left,
  numberstyle=\tiny\color{gray},
  tabsize=4
}

\input{../../Comments}
\input{../../Common}

\begin{document}

\title{Technical Report: Slime Mold Algorithm\\Bio-Inspired Simulation in Ad-Natura\\[0.5em]
\large \progname}
\author{\authname}
\date{\today}

\maketitle

\pagenumbering{roman}

\newpage

\tableofcontents

\newpage

\section{Revision History}

\begin{tabularx}{\textwidth}{p{3cm}p{2cm}p{2.5cm}X}
\toprule {\bf Date} & {\bf Version} & {\bf Author} & {\bf Notes}\\
\midrule
Mar 7, 2026 & 1.0 & Andy Liang & Initial report: architecture, implemented behaviors, design rationale\\
Mar 21, 2026 & 1.1 & Andy Liang & Added behavioral drift, water consumption, metabolic glow; updated design rationale with Sims and Adamatzky\\
Mar 26, 2026 & 2.0 & Andy Liang & Final version: microphone input, consent/resistance, light avoidance implemented; TA feedback addressed; trail deposition corrected; Future Work moved after Conclusion\\
\bottomrule
\end{tabularx}

\newpage

\pagenumbering{arabic}

% ============================================================
\section{Introduction}
\label{sec:intro}
% ============================================================

This report documents the design and implementation of the slime mold
simulation system in \textit{Ad-Natura}, a 2D environmental exploration game.
The simulation models \textit{Physarum polycephalum}, a single-celled organism
capable of decentralized decision-making, adaptive network formation, and
environmental sensing---without any central nervous system or governing
structure \cite{keller1983, adamatzky2014}.

The system runs entirely on the GPU as a compute shader, driving 1024
autonomous agents that deposit chemical trails, sense their environment, and
collectively produce emergent network structures. Beyond serving as a game
mechanic, the simulation is designed as a site of inquiry: a counter-algorithmic
system \cite{bunton2025} that resists optimization in favor of organic
unpredictability.

This report covers the technical architecture (Section~\ref{sec:arch}), agent
behavior (Section~\ref{sec:agents}), trail dynamics
(Section~\ref{sec:trails}), noise generation (Section~\ref{sec:noise}), game
integration (Section~\ref{sec:game}), microphone input
(Section~\ref{sec:mic}), consent and resistance (Section~\ref{sec:consent}),
design rationale (Section~\ref{sec:rationale}), and technical challenges
(Section~\ref{sec:challenges}).

% ============================================================
\section{System Architecture}
\label{sec:arch}
% ============================================================

The simulation is implemented as a Unity compute shader
(\texttt{Slime.compute}) driven by a C\# controller (\texttt{Slime.cs}). It
uses a ping-pong double-buffering scheme and two GPU kernels dispatched each
\texttt{FixedUpdate} frame.

\subsection{Data Structures}

Each agent is represented by a minimal struct:

\begin{lstlisting}
struct Agent {
    float2 position;   // world-space position
    float  angle;      // heading in radians
    float4 type;       // species color (unused in current build)
};
\end{lstlisting}

Agents are stored in a \texttt{StructuredBuffer} of 1024 elements, initialized
with random positions and angles via a uint-based hash PRNG.

\subsection{Ping-Pong Buffering}

Two \texttt{RenderTexture} buffers (\texttt{TrailMap} and
\texttt{TrailMapProcessed}) alternate roles each frame. The \textbf{Update}
kernel reads from the current trail map and writes agent deposits to it. The
\textbf{Postprocess} kernel reads from that same buffer, applies diffusion,
decay, and pulse, then writes the result to the other buffer. On the next
frame, the roles swap. This prevents read-write hazards on the GPU.

Both trail maps use a floating-point texture format
(\texttt{RenderTextureFormat.RFloat} for attraction maps,
\texttt{RenderTextureFormat.ARGBFloat} for trail maps). This is critical:
trail intensity values routinely exceed 1.0, creating a ``reservoir effect''
where dense trail networks store accumulated chemical signal. Clamping to
$[0, 1]$ destroys this behavior entirely.

\subsection{Kernel Overview}

\begin{itemize}
    \item \textbf{Update kernel} (\texttt{[numthreads(16,1,1)]}): One thread
    per agent. Handles sensing, turning, movement, boundary behavior, and trail
    deposition. Dispatched as $\lceil 1024 / 16 \rceil = 64$ thread groups.

    \item \textbf{Postprocess kernel} (\texttt{[numthreads(8,8,1)]}): One
    thread per texel. Handles 3$\times$3 diffusion, entropic decay,
    cytoplasmic pulse, and metabolic glow. Dispatched over the full trail map resolution.
\end{itemize}

% ============================================================
\section{Agent Behavior and Chemotaxis}
\label{sec:agents}
% ============================================================

The simulation models the organism as a population of 1024 \textit{agents}---autonomous computational particles, each representing a small region
of the plasmodial network. An agent has a position, a heading angle, and no
memory beyond these two values. It cannot communicate with other agents
directly; all coordination happens indirectly through chemical trails deposited
in the shared environment (stigmergy). Each agent performs a
sense-turn-move-deposit cycle inspired by the Jones model of Physarum
chemotaxis \cite{adamatzky2014}.

\subsection{Sensing}

The \texttt{sense()} function samples a $3 \times 3$ neighborhood at a point
offset from the agent's position by \texttt{sensorOffsetDist} in a given
direction. It sums both trail intensity and water attraction values, weighted by
\texttt{waterWeight}:

\begin{equation}
    S(\theta) = \sum_{i=-1}^{1}\sum_{j=-1}^{1}
    \Big[\text{trail}(p + d_\theta + (i,j)) +
    w \cdot \text{water}(p + d_\theta + (i,j))\Big]
\end{equation}

where $p$ is the agent position, $d_\theta$ is the sensor offset vector at
angle $\theta$, and $w$ is the water attraction weight.

\subsection{Turning}

Three sensor readings are taken: forward ($S_F$), left ($S_L$, at
$+\texttt{sensorAngleSpacing}$), and right ($S_R$, at
$-\texttt{sensorAngleSpacing}$). The turning logic follows:

\begin{enumerate}
    \item If $S_F > S_L$ and $S_F > S_R$: continue straight.
    \item If $S_F < S_L$ and $S_F < S_R$: turn randomly left or right
    (stochastic tie-breaking prevents deadlock).
    \item If $S_L > S_R$: turn left by \texttt{turnSpeed} $\cdot \Delta t$.
    \item If $S_R > S_L$: turn right by \texttt{turnSpeed} $\cdot \Delta t$.
\end{enumerate}

This produces emergent path-following and network formation without any global
coordination---each agent responds only to local chemical gradients.

\subsection{Movement and Boundary Handling}

After sensing and turning, each agent advances along its heading. Speed is
a continuous blend between an idle base speed and an attracted speed, governed
by local water attraction:

\begin{equation}
    v = \text{lerp}\big(\texttt{idleSpeed},\;
    \texttt{moveSpeed} \cdot \texttt{attractedSpeedMultiplier},\;
    \text{saturate}(\text{attraction})\big)
\end{equation}

This produces smooth acceleration toward water rather than a binary
fast/slow switch. Agents far from any stimulus still drift slowly, keeping the
colony alive even in barren regions.

At boundaries, the behavior depends on whether water is nearby. If the agent
hits an edge where the water attraction map reads above a small threshold, it
deflects \textit{along} the boundary (horizontal or vertical, chosen
randomly) with a random angular spread. This prevents the pile-up that
occurs when agents bounce straight back and immediately re-deposit on the
same pixel. If no water is sensed at the edge, the agent receives a
random heading biased away from the wall, with slight randomness to
avoid artificial symmetry. In both cases, the agent's position is clamped to
stay within bounds.

\subsection{Trail Deposition}

Each agent overwrites its current texel with the organism's base color at a
fixed intensity:

\begin{equation}
    \texttt{TrailMap}[\text{pos}] = \mathbf{c}_{\text{slime}} \cdot
    \text{trailStr}
\end{equation}

where $\text{trailStr} = 1.0$ under normal conditions. The deposit is a hard
overwrite, not an addition---a single agent on an empty pixel produces the
same value as ten agents on the same pixel. Veins emerge not from additive
stacking but from \textit{reinforcement}: many agents repeatedly refreshing
the same path before evaporation can fade it. Unpopular paths evaporate
away; popular paths persist. The \texttt{trailWeight} parameter modulates
how strongly agents \textit{sense} trails (Section~\ref{sec:agents}), not how
much they deposit.

\subsection{Behavioral Drift}
\label{sec:drift}

In a real \textit{Physarum} population, no two regions of the plasmodium behave
identically. To model this, each agent's parameters are modulated by a
per-agent noise seed derived from its index:

\begin{equation}
    p_i = p_{\text{base}} \cdot (1 + m \cdot (2 \cdot \text{hash}(i, t) - 1))
\end{equation}

where $p_i$ is the agent's effective parameter value, $p_{\text{base}}$ is the
shared base value, $m$ is the drift multiplier, and $\text{hash}(i, t)$ is a
slowly evolving noise value seeded by the agent's index. Four parameters are
modulated independently: sensor length ($m = 0.7$), turn speed ($m = 0.6$),
trail weight ($m = 0.8$), and water attraction strength ($m = 0.9$).

The drift evolves over time (controlled by \texttt{driftSpeed}), so an agent
that is currently exploratory may later become conservative, and vice versa.
The population is simultaneously individual \textit{and} collective---what
Sinclair \cite{sinclair2020} describes as paraconsistent individuality, where
an entity can hold contradictory states without resolution. The emergent effect
is a simulation that never fully converges to a single stable pattern.

% ============================================================
\section{Trail Dynamics}
\label{sec:trails}
% ============================================================

The Postprocess kernel transforms the raw trail map each frame through four
sequential stages: diffusion, entropic decay, cytoplasmic pulse (with
brightness tapering), and metabolic glow.

\subsection{Diffusion}

A standard $3 \times 3$ box blur computes the mean trail value in the
neighborhood of each texel. The result is linearly interpolated with the
original value:

\begin{equation}
    T'(x,y) = \text{lerp}\Big(T(x,y),\;\frac{1}{9}\sum_{i,j \in
    \{-1,0,1\}} T(x+i,\; y+j),\;\texttt{diffuseSpeed} \cdot \Delta t\Big)
\end{equation}

This simulates the physical diffusion of chemical signal through the
environment, smoothing the trail network over time.

\subsection{Entropic Decay}
\label{sec:entropy}

Rather than applying a uniform evaporation rate, the system modulates decay
with a two-octave noise field that varies in both space and time:

\begin{equation}
    \text{noise}(x,y,t) = \texttt{organicNoise}\Big(\frac{x \cdot
    s_e}{W},\;\frac{y \cdot s_e}{H},\; t \cdot v_e\Big)
\end{equation}

\begin{equation}
    e(x,y,t) = \texttt{evaporateSpeed} \cdot \Delta t \cdot
    \big(1 + \texttt{entropyStrength} \cdot (2 \cdot \text{noise} - 1)\big)
\end{equation}

\begin{equation}
    T''(x,y) = T'(x,y) - e(x,y,t)
\end{equation}

where $s_e$ is the entropy scale, $v_e$ is the entropy speed, $W$ and $H$ are
texture dimensions. The noise field creates regions of faster and slower decay,
producing organic, non-uniform trail dissolution. Some areas persist longer
while others erode quickly, breaking the visual monotony of uniform evaporation.

This feature is grounded in the observation that real biological decay is never
uniform---it is shaped by local moisture, temperature, microbial activity, and
countless other environmental variables \cite{bunton2025}.

\subsection{Cytoplasmic Streaming (Pulse)}
\label{sec:pulse}

Real \textit{Physarum} exhibits rhythmic cytoplasmic streaming---the
shuttle flow of protoplasm that oscillates through the organism's network, driving
nutrient transport and enabling distributed decision-making \cite{adamatzky2014}.

The simulation reproduces this as a density-thresholded brightness oscillation:

\begin{equation}
    d_n = \text{smoothstep}(0.3,\; 0.6,\; T''(x,y) / T_{\max})
\end{equation}

\begin{equation}
    f(x,y) = \text{lerp}\big(\texttt{pulseMinFreq},\;
    \texttt{pulseMaxFreq},\; \text{water}(x,y)\big)
\end{equation}

\begin{equation}
    \phi(x,y) = (x + y) \cdot \texttt{pulseWaveScale}
\end{equation}

\begin{equation}
    \text{pulse}(x,y,t) = d_n \cdot \texttt{pulseAmplitude} \cdot
    \sin(2\pi f t + \phi)
\end{equation}

The final trail output is:

\begin{equation}
    T_{\text{out}}(x,y) = \max\big(T''(x,y) + \text{pulse}(x,y,t),\; 0\big)
\end{equation}

Key design decisions:
\begin{itemize}
    \item \textbf{Density threshold}: The smoothstep activation ensures only
    established trail networks pulse, not sparse background noise. This mirrors
    real Physarum, where streaming occurs in mature veins.

    \item \textbf{Water-driven frequency}: Pulse frequency increases near water
    sources, simulating heightened metabolic activity near nutrients.

    \item \textbf{Spatial phase offset}: The $(x+y)$ phase term creates the
    illusion of wave propagation across the network, rather than uniform
    blinking.

    \item \textbf{No clamping}: The output uses $\max(\cdot, 0)$ rather than
    $\text{clamp}(0, 1)$, preserving the reservoir effect of accumulated
    chemical signal.

    \item \textbf{Brightness taper}: Because the pulse operates within a
    ping-pong buffer, its multiplicative boost can ratchet brightness upward
    across frames in dense clusters. A soft taper scales the pulse contribution
    down as brightness approaches 2.0:
    $\text{taper} = \text{saturate}(2.0 - \|\mathbf{T}\|)$. This preserves the
    glow in dense networks while preventing blowout.
\end{itemize}

\subsection{Metabolic Glow}
\label{sec:glow}

When the slime mold consumes water from the environment
(Section~\ref{sec:consumption}), a \texttt{metabolicGlow} value is tracked on
the CPU and passed to the shader. The Postprocess kernel applies this as an
additive brightness boost, localized to regions where trails and water overlap:

\begin{equation}
    T_{\text{out}}(x,y) = T_{\text{pulse}}(x,y) +
    \mathbf{c}_{\text{slime}} \cdot g \cdot 0.05 \cdot d_n \cdot w(x,y)
\end{equation}

where $g$ is the current metabolic glow intensity,
$\mathbf{c}_{\text{slime}}$ is the organism's base color, $d_n$ is the
same density threshold used by the pulse, and $w(x,y)$ is the local water
attraction value. This localizes the glow to feeding sites rather than
illuminating the entire network---the bloom appears where the organism meets
the water, not everywhere it has been. The glow is strictly additive, not
multiplicative---a critical design decision. Because the Postprocess kernel
operates within a ping-pong buffer, any multiplicative amplification
($T \gets T \cdot (1 + g)$) would compound exponentially across frames,
producing catastrophic feedback. Additive glow reaches a natural equilibrium
with evaporation, producing a visible but stable brightness increase that
decays smoothly as the glow value fades on the CPU side.

The effect is a brief luminous bloom when the organism feeds---a visible
metabolic event that makes the act of consumption perceptible.

% ============================================================
\section{GPU Noise Implementation}
\label{sec:noise}
% ============================================================

GPU-based noise requires careful implementation. Standard approaches using
$\sin()$-based hash functions suffer from precision artifacts on many GPU
architectures. The system uses an integer hash approach:

\begin{lstlisting}
float hashNoise(float2 p) {
    uint2 i = uint2(abs(p));
    uint  h = i.x * 1597334677u ^ i.y * 3812015801u;
    h = h * h;
    return float(h) / 4294967295.0;
}
\end{lstlisting}

This provides well-distributed pseudorandom values without relying on
trigonometric functions. A \texttt{valueNoise()} function interpolates between
hash corners using smoothstep, and \texttt{organicNoise()} layers two octaves
with halved amplitude:

\begin{equation}
    \text{organicNoise}(x, y, t) = \text{valueNoise}(x, y, t) +
    0.5 \cdot \text{valueNoise}(2x, 2y, 2t)
\end{equation}

The result is normalized to $[0, 1]$ by dividing by 1.5. This two-octave
structure adds fine detail to the noise field without the computational cost of
full fractal Brownian motion, which is unnecessary for the visual effect.

% ============================================================
\section{Game Integration}
\label{sec:game}
% ============================================================

While the simulation is designed as a self-contained system, it integrates with
the broader game through several mechanics:

\begin{itemize}
    \item \textbf{Water Attraction}: A water attraction map
    (\texttt{WaterAttractionMap}) is generated by the game's water simulation
    system and passed to the compute shader. Agents preferentially move toward
    water, deposit stronger trails near it, and pulse at higher frequencies in
    its proximity.

    \item \textbf{Decomposition}: The slime mold participates in the game's
    decomposition system, breaking down organic obstacles in the environment.
    Trail density in an area contributes to the decomposition rate of nearby
    objects.

    \item \textbf{Water Consumption}: \label{sec:consumption} The organism
    actively consumes water from the game's cellular liquid simulation. The C\#
    controller periodically samples trail density across the texture, identifies
    regions where both slime trails and water are present, and calls
    \texttt{RemoveWater()} on the liquid simulation to drain a cluster of
    neighboring grid cells. A \texttt{metabolicGlow} value accumulates during
    consumption and decays over time, driving the visual glow in the shader
    (Section~\ref{sec:glow}). This creates a reciprocal material
    relationship: the water shapes the organism's growth, and the organism
    depletes the water. Neither remains unchanged by the encounter---what Barad
    \cite{barad2011} terms intra-action, where entities co-constitute through
    mutual transformation rather than one-directional influence.

    \item \textbf{Light Avoidance}: Real \textit{Physarum} is negatively
    phototactic. The simulation encodes this via the R channel of the
    \texttt{WaterAttractionMap}, which stores a light fear value. Agents in lit
    areas blend between trail-following chemotaxis and light-fleeing behavior:
    fear dominates but trail signal still contributes, so the organism
    \textit{resists} light rather than losing all self-organization. This
    preserves emergent network behavior under environmental pressure.

    \item \textbf{Hazard System}: Calamity objects (encoded in a separate
    \texttt{HazardMap}) corrode trails, slow agents, and kill agents on direct
    contact. Dead agents leave a colored death mark at their last position and
    respawn at the colony center. The stain falloff is pre-computed on the CPU
    to avoid expensive per-pixel GPU neighbor searches.
\end{itemize}

The \texttt{SlimeMoldManager} class coordinates between the simulation and
other game systems, providing a public API (\texttt{SetAttractionMaps},
\texttt{GetTrailTexture}, \texttt{GetWorldBounds}) for other scripts to
interact with the slime mold.

% ============================================================
\section{Design Rationale}
\label{sec:rationale}
% ============================================================

The design of this system is informed by both biological research and critical
theory. The artist's role here is not that of a central author imposing form
onto matter, but a designer of conditions under which matter performs
\cite{adamatzky2016}. Several deliberate choices distinguish it from
conventional algorithmic implementations.

\subsection{Self-Organization Without Central Control}

Evelyn Fox Keller's foundational work on slime mold aggregation \cite{keller1983}
demonstrated that scientists historically assumed \textit{Dictyostelium}
required a ``pacemaker'' cell to coordinate collective behavior. Keller showed
that aggregation is fully self-organizing: no leader, no central governor.

This principle is encoded in the simulation's architecture. All 1024 agents
follow identical local rules---sense, turn, move, deposit. Network formation,
trail reinforcement, and path optimization emerge entirely from these local
interactions. There is no global pathfinding, no master controller, no
optimization objective.

\subsection{Externalized Memory}

The trail map functions as a form of externalized spatial memory. As Sims
\cite{sims2021} observes, slime mold does not store memory internally; it
remembers through the chemical traces it leaves in its environment. In the
simulation, the trail network is simultaneously the organism's history, its
communication medium, and the substrate for collective decision-making.
Entropic decay (Section~\ref{sec:entropy}) ensures this memory is impermanent
and shaped by environmental noise---not a perfect record but an organic,
degrading trace.

\subsection{Counter-Algorithmic Design}

Bunton \cite{bunton2025} critiques the Slime Mold Algorithm (SMA) for reducing
\textit{Physarum}'s complex, ``queer'' biology to a metaheuristic optimization
tool---stripping away precisely the behaviors that make the organism
interesting. The concept of ``jamming''---blocking, improvising, knotting---is
proposed as an alternative relation to algorithmic systems.

The entropic decay system (Section~\ref{sec:entropy}) functions as a form of
jamming. Rather than optimizing trail networks toward efficiency, the
noise-modulated decay introduces unpredictable dissolution. Trails that
``should'' persist may erode; regions that ``should'' clear may linger. The
system resists convergence toward optimal solutions.

\subsection{Nature's Performativity}

Barad's theory of agential realism \cite{barad2011} argues that nature is not
passive matter awaiting measurement, but actively performs---entities
co-constitute through ``intra-action'' rather than interaction. The slime mold
simulation embodies this: the organism and its environment are not separate
systems that interact, but a single entangled process. Trail deposition shapes
agent movement, which shapes trail deposition. The water map shapes the
organism, and in the game's decomposition system, the organism reshapes its
environment.

\subsection{Non-Human Play}

Stone \cite{stone2019} extends the concept of ``player'' beyond the human,
drawing on Bennett's distributive agency to argue that nonhuman entities---bacteria, plants, algorithms---can be understood as players in an
assemblage. In \textit{Ad-Natura}, the slime mold is not an object the player
controls but a co-player with its own agency: it responds to the environment,
makes decisions (through chemotaxis), and produces outcomes the human player
cannot fully predict or control.

This framing connects to Ruberg's concept of ``after agency''
\cite{ruberg2022}---game systems where the human player's control is
deliberately limited. The player can influence the slime mold's environment
(e.g., directing water flow), but cannot command the organism directly.

% ============================================================
\section{Technical Challenges and Lessons}
\label{sec:challenges}
% ============================================================

\subsection{The Reservoir Effect}

The most consequential technical discovery during development was that trail
values in floating-point \texttt{RenderTexture} buffers routinely exceed 1.0.
When multiple agents deposit trails in the same region over several frames, the
accumulated value grows well above the $[0, 1]$ range. This ``reservoir
effect'' is essential to the simulation's visual quality: dense networks appear
brighter and more saturated, while sparse trails fade naturally.

An early implementation clamped the output to $[0, 1]$, which flattened all
dense networks to a uniform maximum brightness and eliminated the visual
distinction between mature veins and new trails. The fix was simple but
non-obvious: output $\max(\text{value}, 0)$ with no upper clamp.

\subsection{GPU Noise Precision}

Initial noise implementations using $\sin()$-based hash functions produced
visible banding artifacts on certain GPU architectures. The integer hash
approach (Section~\ref{sec:noise}) resolved this entirely, at no performance
cost.

\subsection{RenderTexture Format}

The default water attraction map must use
\texttt{RenderTextureFormat.RFloat}, not \texttt{ARGB32}. Using an integer
format caused the attraction values to be silently quantized, producing
step-like artifacts in agent movement near water boundaries.

\subsection{Inspector Serialization}

Unity's Inspector serializes field values independently of code defaults. When
default parameter values were changed in the C\# script, existing scene
references retained their old serialized values. This caused confusion during
development, as code changes appeared to have no effect until the component was
reset manually.

% ============================================================
\section{Microphone as Interface}
\label{sec:mic}
% ============================================================

A standalone \texttt{MicrophoneInput} component captures audio from the
system's default recording device via Unity's \texttt{Microphone} API. Three
features are extracted each frame:

\begin{itemize}
    \item \textbf{Amplitude}: RMS energy of the most recent 1024 samples,
    scaled by a configurable gain and smoothed via exponential lerp.
    \item \textbf{Spectral centroid}: The energy-weighted mean frequency bin,
    computed from Unity's \texttt{GetSpectrumData} FFT (Blackman-Harris
    window). Normalized to $[0, 1]$, where 0 is low-pitched and 1 is
    high-pitched.
    \item \textbf{Onset}: The positive difference in amplitude between
    consecutive frames, thresholded and scaled to detect sudden sound events.
\end{itemize}

A hidden \texttt{AudioSource} (muted, zero volume) plays the microphone clip
to provide Unity's FFT pipeline with data. The source is muted to prevent
audible feedback; \texttt{GetSpectrumData} still functions when muted.

\subsection{Parameter Mapping}

The C\# controller reads mic values into local copies of simulation parameters
each frame, leaving Inspector-serialized values untouched. A master
\texttt{micInfluence} slider scales all effects:

\begin{itemize}
    \item \textbf{Silence} ($\text{amplitude} \approx 0$): Pulse frequency
    lifts gently. The organism notices the player's quiet presence.
    \item \textbf{Loud sound}: Evaporation rate increases (trails thin),
    turn speed increases (agents agitate). The network becomes restless.
    \item \textbf{High pitch}: Sensor range contracts, turn speed increases
    further. The organism becomes nervous, myopic.
    \item \textbf{Low pitch}: Diffusion rate increases. The trail network
    softens and spreads.
    \item \textbf{Sudden onset}: Trail sensing weight spikes briefly---a
    short-term memory of the sound event.
\end{itemize}

Turn speed modulation from amplitude and pitch are combined additively
(not multiplicatively) to prevent compounding that spirals agents into dense
clusters. Sensor shrink is floored at 50\% of base length to ensure agents
always retain enough sensing range to disperse.

% ============================================================
\section{Consent and Resistance}
\label{sec:consent}
% ============================================================

The organism can refuse the player's input. A stress accumulator on the CPU
tracks how long the microphone amplitude exceeds a configurable threshold.
When stress crosses a second threshold, the organism enters a
\textit{refractory state}---a timed withdrawal during which it becomes
sluggish, dim, and unresponsive.

\subsection{Stress Dynamics}

\begin{equation}
    \sigma_{t+1} = \begin{cases}
    \sigma_t + (\text{amp} - \theta_s) \cdot r_a \cdot \Delta t
        & \text{if } \text{amp} > \theta_s \\
    \sigma_t - r_d \cdot \Delta t & \text{otherwise}
    \end{cases}
\end{equation}

where $\sigma$ is the stress level, $\theta_s$ is the stress threshold,
$r_a$ is the accumulation rate, and $r_d$ is the recovery rate. The transition
into withdrawal is not instantaneous: as stress builds, the organism's
responsiveness to mic input fades proportionally
($\text{responsiveness} = 1 - 0.7\sigma$), so the player can observe the
organism becoming less reactive before it shuts down entirely.

\subsection{Refractory State}

When $\sigma \geq \theta_r$ (the refractory threshold), the organism enters
withdrawal. A smoothed refractory signal ramps up gradually (preventing a
harsh visual snap) and is passed to the GPU, where it dampens:

\begin{itemize}
    \item \textbf{Behavioral drift}: individuality collapses, the collective
    becomes uniform ($\times 0.9$ reduction).
    \item \textbf{Speed}: the organism becomes sluggish ($\times 0.35$
    reduction).
    \item \textbf{Trail deposit}: trails dim ($\times 0.2$ reduction).
    \item \textbf{Pulse amplitude}: breathing shallows ($\times 0.4$
    reduction).
\end{itemize}

Recovery requires silence. If the microphone amplitude remains above the
stress threshold during withdrawal, the recovery timer pauses---the organism
cannot heal while still being assaulted. This is a deliberate design
choice grounded in Ruberg's ``after agency'' \cite{ruberg2022}: the player's
control is not removed by a game mechanic but rendered meaningless by the
organism's refusal.

% ============================================================
\section{Conclusion}
\label{sec:conclusion}
% ============================================================

The slime mold simulation in \textit{Ad-Natura} is a GPU-accelerated,
agent-based model of \textit{Physarum polycephalum} that prioritizes biological
plausibility and emergent complexity over algorithmic efficiency. Through
entropic decay, cytoplasmic pulse, behavioral drift, water consumption,
metabolic glow, real-time microphone input, and consent-based withdrawal, the
system produces organic, unpredictable behavior from simple local
rules---while engaging in reciprocal exchange with both its digital
environment and the player's physical body.

The microphone transforms the simulation from a reactive display into a site
of encounter: the player's breath, voice, and silence all register as input,
and the organism responds on its own terms---including the capacity to refuse.
The design is grounded in critical engagement with biological research and
theoretical frameworks that challenge conventional approaches to computational
modeling of living systems. Rather than reducing \textit{Physarum} to an
optimization heuristic, the system honors the organism's capacity for
self-organization, distributed agency, externalized memory, and resistance to
centralized control.

% ============================================================
\begin{thebibliography}{9}

\bibitem[Adamatzky(2014)]{adamatzky2014}
A.~Adamatzky.
\newblock Slime mould analog models of space exploration.
\newblock \textit{International Journal of Bifurcation and Chaos}, 24(02):1430006, 2014.

\bibitem[Adamatzky and Schubert(2016)]{adamatzky2016}
A.~Adamatzky and T.~Schubert.
\newblock \textit{Experiencing the Unconventional: Science in Art}.
\newblock World Scientific, 2016.

\bibitem[Sims(2021)]{sims2021}
B.~Sims.
\newblock \textit{Slime Mould and Philosophy}.
\newblock Punctum Books, 2021.

\bibitem[Barad(2011)]{barad2011}
K.~Barad.
\newblock Nature's queer performativity.
\newblock \textit{Qui Parle}, 19(2):121--158, 2011.

\bibitem[Bunton(2025)]{bunton2025}
K.~Bunton.
\newblock \textit{Jamming Playbook: Re/Decomposing the Slime Mold Algorithm}.
\newblock MA thesis, Royal Academy of Art, The Hague, 2025.

\bibitem[Keller(1983)]{keller1983}
E.~F.~Keller.
\newblock The force of the pacemaker concept in theories of aggregation.
\newblock \textit{Journal of the History of Biology}, 16(1):39--56, 1983.

\bibitem[Ruberg(2022)]{ruberg2022}
B.~Ruberg.
\newblock After agency: Video games and the non-player character.
\newblock In \textit{The Routledge Companion to Media and Fairy-Tale Cultures}, pages 479--487. Routledge, 2022.

\bibitem[Sinclair(2020)]{sinclair2020}
N.~W.~Sinclair.
\newblock Exploding individuals: Engaging indigenous logic within science.
\newblock \textit{Hypatia}, 35(1):58--74, 2020.

\bibitem[Stone(2019)]{stone2019}
K.~Stone.
\newblock What can play: The potential of non-human players.
\newblock In \textit{Games and Play in the Creative, Media and Cultural Industries}, pages 165--180. Routledge, 2019.

\end{thebibliography}

\end{document}
